\section{Functions}
\subsection{Compositions and Inverses}
\subsubsection{Compositions}
\begin{definition} \label{def:f-composition}
If $f$ and $g$ are two functions, we can define a new function $fg$ by 
\[ fg(x) = f(g(x)) \]
Note that this relies on $f$ being defined at the point $g(x)$. If $f$ has a \emph{singularity} at a point $a$ (as in $f(a)$ divides by zero somewhere in the calculation) then if there is an $x$ such that $g(x) = a$, $fg(x) = f(a)$, which would not be allowed.
\end{definition}

We can think of this similar to bracket rules - you do the ``innermost'' one first, and work outwards. 
In the case of Definition \ref{def:f-composition} we compute $g(x)$ first, and then $f(x)$

This can keep going! If we have three functions $f$, $g$, $h$, then $fgh(x) = f(g(h(x)))$, and we evaluate the functions inside-out ($h$ acts on $x$ first, then $g$, then $f$).
\subsubsection{Compositions Examples and Exercises}
\begin{example} \label{ex:simple-composition}
    Let $f$ be a function defined by 
    \[ f(x) = \frac{2x + 1}{5}\]
    Find $ff(x)$, and write your answer in the form
    \[ff(x) = \frac{ax + b}{25}\]
    Where $a,b$ are numbers.
\end{example}
\begin{solution}
    First write:
    \[ f(x) = \frac{2x + 1}{5}\]
    So 
    \[ ff(x) = f(f(x)) = \frac{2\textcolor{red}{f(x)} + 1}{5}\]
    Substituting $f(x)$ into this equation gives 
    \[ ff(x) =  \frac{2\frac{2x + 1}{5} + 1}{5}\]
    Here we have a fraction with fractions in it. Fractions are annoying, so we want to reduce the number of fractions. Notice that if we times the inner fraction by $5$, it would remove the fraction. 
    We can use the fact that any number times $1$ is itself to multiply the fraction by $5/5$:
    \begin{align*}
        ff(x) &=  \frac{5}{5} \times \frac{2\frac{2x + 1}{5} + 1}{5} & \tag*{(substitute in $f(x)$)}\\
            &= \frac{2\times \cancel{5} \times \frac{2x + 1}{\cancel{5}} +5}{5 \times 5} & \tag*{(expand top and bottom)}\\
            &= \frac{4x + 2 + 5}{25} \\
            &= \frac{4x + 7}{25}
    \end{align*}

    Now set 
    \[ \frac{4x + 7}{25} = \frac{ax + b}{25}\]
    Which finishes the question, finding $a = 4$, $b = 7$.
    \qed
\end{solution}



\begin{example} \label{ex:simple-composition2}
    Let $f$ be a function defined by 
    \[ f(x) = \frac{2x + 1}{4x}\]
    Where $x \neq 0$. \emph{This is required to not divide by $0$.}

    Find $ff(x)$, and write your answer in the most simplified form of
    \[ff(x) = \frac{ax + b}{cx + d}\]
    Where $a,b,c,d$ are numbers. Make sure to state where $ff(x)$ is defined (what can $x$ not be).
\end{example}
\begin{solution}
    First write (whenever $x\neq 0$):
   \[ f(x) = \frac{2x + 1}{4x}\]
    So 
    \[ ff(x) = f(f(x)) = \frac{2\textcolor{red}{f(x)} + 1}{4\textcolor{red}{f(x)}}\]
    Substituting $f(x)$ into this equation gives 
    \[ ff(x) =  \frac{2\frac{2x + 1}{4x}\ + 1}{4 \frac{2x + 1}{4x}}\]
    Here we have a fraction with fractions in it. Fractions are annoying, so we want to reduce the number of fractions. Notice that if we times the inner fractions by 
    $4x$, it would remove the fraction. 
    We can use the fact that any number times $1$ is itself to multiply the fraction by $4x/4x$, which is allowed because $x\neq 0$:
    \begin{align*}
        ff(x) &=  \frac{4x}{4x} \times \frac{2\frac{2x + 1}{4x}\ + 1}{4 \frac{2x + 1}{4x}} & \tag*{(substitute in $f(x)$)}\\
            &= \frac{2\times \cancel{4x} \times \frac{2x + 1}{\cancel{4x}} + 4x}{4 \times \cancel{4x} \times \frac{2x + 1}{\cancel{4x}}} & \tag*{(expand top and bottom)}\\
            &= \frac{4x + 2 + 4x}{8x + 4} \\
            &= \frac{8x + 2}{8x + 4} \\
            &= \frac{2}{2} \times \frac{4x + 1}{4x + 2} & \tag*{(factorise top and bottom)}\\
            &= \frac{4x + 1}{4x + 2}
    \end{align*}

    Since we cannot simplify the equation any further, we now set 
    \[ \frac{4x + 1}{4x + 2} = \frac{ax + b}{cx + d}\]
    Which finds $a = c = 4$, $b = 1$, $d = 2$ by comparison.
    $ff(x)$ cannot divide by zero, so we know  $4x+2 \neq 0$, and get that therefore  
    $x \neq - \frac{1}{2} $. But remember $f$ itself requires $x \neq 0$, so $x$ cannot be $0$ or $-\frac{1}{2}$.
    Which finishes the question.
    \qed 
\end{solution}

\begin{example} \label{ex:simple-composition-different-functions}
    Let $f$ and $g$ be functions defined by 
    \[ f(x) = \sin (x) + \cos (3x) \]
    \[g(x) = \frac{x +1}{3}\]

    Find $fg(x)$, simplifying your answer.
\end{example}
\begin{solution}
    First write:
    \[ f(x) = \sin (x) + \cos (3x) \]
    \[g(x) = \frac{x +1}{3}\]
    So 
    \[ fg(x) = f(g(x)) = \sin (\textcolor{red}{g(x)}) + \cos (3\textcolor{red}{g(x)})\]
    Substituting $g(x)$ into this equation gives 
    \begin{align*} fg(x) &=  \sin \left(\frac{x +1}{3}\right) + \cos \left(3\times\frac{x +1}{3}\right) \\&= \sin \left(\frac{x +1}{3}\right) + \cos (x +1)
    \end{align*}
    Which finishes the question.
    \qed
\end{solution}
\begin{exercise} \label{ex:composition-three-functions}
    \emph{Try it yourself!}
    
    Let $f$, $g$, $h$, be functions defined by 
    \[ f(x) = \frac{1}{2}x^2\]
    \[g(x) = \sin(x) - \cos(x)\]
    \[h(x) = \cos(x) + \sin(x)\]

    Show that \[ fg(x) + fh(x) = 1\] 
    Hint: you may use (or recall) without proof that $\sin(x)^2 + \cos(x)^2 = 1$.
\end{exercise}
\begin{solution}
    We will approach this question by calculating $fg$ and $fh$ separately, and then plugging them in.
    First write:
    \[ f(x) = \frac{1}{2}x^2\]
    \[g(x) = \sin(x) - \cos(x)\]
    \[h(x) = \cos(x) + \sin(x)\]
    So 
    \[ fg(x) = f(g(x)) =\frac{1}{2}\textcolor{red}{g(x)}^2\]
    Substituting $g(x)$ into this equation gives 
    \[ fg(x) = \frac{1}{2}(\sin(x) - \cos(x))^2 \]

    Similarly, 
    \[ fh(x) = f(h(x)) = \frac{1}{2}\textcolor{red}{h(x)}^2 \]
    Subsituting $h(x)$ into this equation gives
    \[ fh(x) = \frac{1}{2}(\cos(x) + \sin(x))^2 \]

    Then we can calculate:
    \begin{align*}
        fg(x) + fh(x) &= \frac{1}{2}(\sin(x) - \cos(x))^2 +  \frac{1}{2}(\cos(x) + \sin(x))^2 & \tag{plugging in the equations}\\
        &= \frac{1}{2}\left((\sin(x) - \cos(x))^2 + (\cos(x) + \sin(x))^2\right) & \tag{factoring}\\
        &= \frac{1}{2}\left( \sin(x)^2 - \cancel{2\sin(x)\cos(x)} + \cos(x)^2  + \cos(x)^2 + \cancel{2\sin(x)\cos(x)} + \sin(x)^2 \right) & \tag{expanding} \\
        &= \frac{1}{\cancel{2}} \times \cancel{2} (\sin(x)^2 + \cos(x)^2) & \tag{factoring} \\
        &= \sin(x)^2 + \cos(x)^2 \\
        &= 1 \tag{trig rules (the hint)}
    \end{align*}
    Which finishes the question.
    \qed
\end{solution}

\subsubsection{Inverses} \label{sec:function-inverse}
\begin{definition} \label{def:f-inverse}
    If $f$ is a function, we define the inverse of $f$, denoted $f^{-1}$ by the following equation:
    \[ ff^{-1}(x) = x \]
    Or equivalently, 
    \[ f^{-1}f(x) = x \]
    The most common way to solve for $f^{-1}$ is by: 
    \begin{enumerate} 
        \item Write $y = f(x)$
        \item Rearrange so that we have $x$ in terms of $y$
        \item Swap $x$ and $y$ to find $f^{-1}(x)$
    \end{enumerate}
    You can show that this satisfies the definition of the inverse of $f$. 
    Intuitively, Step 2 above finds ``what do you need to do to $y$ to find $x$''? This is because we are writing $x$ as a bunch of operations 
    applied to $y$, but $y$ was chosen to be $f(x)$ in step $1$, therefore it must ``undo'' the operations we did in Step 1. This argument can be 
    made rigorous.
\end{definition}
Here we have used the notation for function composition in Definition \ref{def:f-composition}.

\subsubsection{Inverses Examples and Exercises}
\begin{example} \label{ex:simple-inverse}
    Example \ref{ex:simple-composition} continued.

    Let $f$ be a function defined by 
    \[ f(x) = \frac{2x + 1}{5}\]

    Find $f^{-1}(x)$, and write your answer in the form
    \[f^{-1}(x) = \frac{ax + b}{2}\]
    Where $a,b$ are numbers.


\end{example}
\begin{solution}
    First write 
    \[y = \frac{2x + 1}{5}\]

    Then rearrange for $x$:
    \begin{align*}
        5y &= 2x+1 \tag{multiply by 5}\\
        5y -1 &= 2x \tag{minus 1}\\
        x &= \frac{5y - 1}{2} \tag{divide by 2}
    \end{align*}
    Therefore,
    \[ f^{-1}(x) = \frac{5x - 1}{2}\]
    And setting it equal to $\frac{ax + b}{2}$, we find that $a = 5$ and $b = -1$ by comparison.
    \qed
\end{solution}

\begin{example} \label{ex:simple-inverse-verification}
    Example \ref{ex:simple-inverse} continued.

    Let $f$ and $g$ be functions defined by 
    \[ f(x) = \frac{2x + 1}{5}\]
    \[ g(x) = \frac{5x - 1}{2}\]
    By computing $fg(x)$, verify that $g$ is the inverse function of $f$, verifying the calculation in the previous question.

\end{example}
\begin{solution}
    First write 
    \[ f(x) = \frac{2x + 1}{5}\]
    \[ g(x) = \frac{5x - 1}{2}\]

    Then calculate
    \begin{align*}
        f(g(x)) &= \frac{2\textcolor{red}{g(x)} + 1}{5} \\
         &= \frac{\cancel{2}\frac{5x - 1}{\cancel{2}} + 1}{5} \\
         &= \frac{5x \cancel{- 1 + 1}}{5}  \\
         &= \frac{\cancel{5}x}{\cancel{5}} \\
         &= x
    \end{align*}
    Therefore, $fg(x) = x$, which verifies $g$ is the inverse function of $f$.
    \qed
\end{solution}



\begin{example} \label{ex:simple-inverse2}
    Example \ref{ex:simple-composition2} continued.

    Let $f$ be a function defined by 
    \[ f(x) = \frac{2x + 1}{4x}\]
    Where $x \neq 0$. \emph{This is required to not divide by $0$.}

    Find $f^{-1}(x)$, and write your answer in the form of
    \[f^{-1}(x) = \frac{1}{ax + b}\]
    Where $a,b$ are numbers. Make sure to state where $f^{-1}(x)$ is defined (what can $x$ not be).


\end{example}
\begin{solution}
    Let $x \neq 0$. First write 
    \[y = \frac{2x + 1}{4x}\]

    Then rearrange for $x$:
    \begin{align*}
        4xy &= 2x+1 \tag{multiply by $4x$}\\
        4xy -2x &= 1 \tag{minus $2x$ to get $x$ on one side}\\
        x(4y - 2) &= 1 \tag{factorise $x$}\\
        x &= \frac{1}{4y-2} \tag{divide by $4y - 2$}
    \end{align*}
    Therefore, 
    \[ f^{-1}(x) = \frac{1}{4x-2}\]
    And setting it equal to $\frac{1}{ax + b}$, we find that $a = 4$ and $b = -2$ by comparison.
    This function is valid (not dividing by zero) whenever $4x- 2 \neq 0$, therefore it is valid whenever $x \neq \frac{1}{2}$.
    \qed
\end{solution}

\begin{exercise} \label{ex:simple-inverse-verification2}
    Example \ref{ex:simple-inverse2} continued.
    \emph{Try it yourself!}

    Let $f$ and $g$ be functions defined by 
    \[ f(x) = \frac{2x + 1}{4x}, \; x \neq 0\]
    \[ g(x) = \frac{1}{4x-2}, \; x \neq \frac{1}{2}\]
    By computing $fg(x)$, verify that $g$ is the inverse function of $f$, verifying the calculation in the previous question.
\end{exercise}
\begin{solution}
    First write 
    \[ f(x) = \frac{2x + 1}{4x}\]
    \[ g(x) = \frac{1}{4x-2} = \frac{1}{2} \times \frac{1}{2x-1}\]

    Then set $x \neq  \frac{1}{2}$ and calculate
    \begin{align*}
        f(g(x)) &= \frac{2\textcolor{red}{g(x)} + 1}{4\textcolor{red}{g(x)}} \\
         &= \frac{\cancel{2}\times\frac{1}{\cancel{2}}\times\frac{1}{2x-1} + 1}{4 \times \frac{1}{2} \times \frac{1}{2x-1}} \\
         &= \frac{\frac{1}{2x-1} + 1}{2 \times \frac{1}{2x-1}}  
    \end{align*}

    Now we use my favourite trick of getting-rid-of-the-annoying-fractions-by-multiplying-by-the-inner-denominator, which is allowed because $x \neq  \frac{1}{2}$:
    \begin{align*}
        f(g(x)) &= \frac{2x-1}{2x-1} \times \frac{\frac{1}{2x-1} + 1}{2 \times \frac{1}{2x-1}}  \tag{multiply by $1$}\\
         &= \frac{\cancel{(2x-1)} \times \frac{1}{\cancel{2x-1}} +2x - 1}{2 \times \cancel{(2x -1) } \times \frac{1}{\cancel{2x-1}}}  \tag{simplify fractions} \\ 
         &= \frac{\cancel{1} + 2x - \cancel{1}}{2} \\
         &= \frac{2x}{2} \\
         &= x
    \end{align*}
    Therefore, $fg(x) = x$, which verifies $g$ is the inverse function of $f$.
    \qed
\end{solution}

\subsubsection{Additional Exercises}

\begin{exercise}
    The function $f$ is defined by 
    \[f(x) = \frac{4x + 3}{x-2}, \;\;\;\;x \neq 2\]
\begin{enumerate}
    \item Find $f^{-1}$
    \item Show that 
    \[ff(x) = \frac{ax + b}{cx + d}\]
    where $a$, $b$, $c$, $d$ are integers to be found.
\end{enumerate}

The point $P$ $(3,15)$ lies on the curve with equation $y = f(x)$.
\begin{enumerate}
\setcounter{enumi}{2}
    \item Find the point to which $P$ is mapped when $y = f(x)$ is transformed to the curve with equation $y = 2f(3x)+8$.
    \item (Bonus) Verify $f^{-1}$ in the first part satisfies $ff^{-1}(x) = x$.
\end{enumerate}
\end{exercise}
\begin{solution}


    \begin{enumerate}
        \item
        Find $f^{-1}$

        First write $y = f(x)$, where $x \neq 2$ i.e.
        \[ y = \frac{4x + 3}{x-2} \]

        Now rearrange for $x$, getting rid of any fractions.
        \begin{align*}
            y &= \frac{4x + 3}{x-2} \\
             y(x-2) &= 4x + 3 \tag{multiply by $(x-2)$} \\
             xy - 2y &= 4x + 3 \tag{expand} \\
             xy - 4x &= 2y+3 \tag{collect $x$ terms to one side} \\
             x(y-4) &= 2y + 3 \tag{factorise LHS} \\
             x &= \frac{2y+3}{y-4} \tag{assuming $y \neq 4$}
        \end{align*}
        Hence, whenever $x \neq 4$, 
        \[ f^{-1}(x) = \frac{2x+3}{x-4}\]
        \item 
        Show  
        \[ff(x) = \frac{ax + b}{cx + d}\]

        Then calculate 
        \begin{align*}
            ff(x) &= f(\textcolor{red}{f(x)}) \\
             &= \frac{4\textcolor{red}{f(x)} + 3}{\textcolor{red}{f(x)}-2}  \tag{insert $f(x)$ wherever $x$ is}\\
             &=  \frac{4\times \textcolor{red}{\frac{4x + 3}{x-2}} + 3}{\textcolor{red}{\frac{4x + 3}{x-2}}-2} \tag{substitute $f(x)$}\\
             &= \frac{x-2}{x-2} \times \frac{4\times \textcolor{red}{\frac{4x + 3}{x-2}} + 3}{\textcolor{red}{\frac{4x + 3}{x-2}}-2} \tag{multiply by $1$ to remove the $x-2$ denominators}\\
             &= \frac{(x-2) (4\times \textcolor{red}{\frac{4x + 3}{x-2}} + 3)}{(x-2)(\textcolor{red}{\frac{4x + 3}{x-2}}-2)} \tag{multiply fractions} \\
             &= \frac{4(4x + 3) + 3(x-2)}{4x + 3 - 2(x-2)} \tag{expand multiplication} \\
             &= \frac{16x + 12 + 3x - 6}{4x + 3 - 2x + 4} \tag{expand multiplication} \\
             &= \frac{19x + 6}{2x + 7} \tag{collect like terms}
        \end{align*}
        This is valid whenever $2x+7 \neq 0$, which is when $x \neq -\frac{7}{2}$

        Hence we have found
        \[ ff(x) = \frac{19x + 6}{2x + 7} \;\;\;\; x \neq 2, \; x \neq -\frac{7}{2}\]

        \item 
        We will reword the question slightly. Define $g$ by 
        \[g(x) = 2f(3x) + 8\]
        Find $g(3)$.

        First calculate $f(\textcolor{red}{3x})$.
        \[ f(\textcolor{red}{3x}) = \frac{4\times \textcolor{red}{3x} + 3}{\textcolor{red}{3x}-2} = \frac{12x + 3}{3x-2}\]
        
        Therefore 

        \begin{align*}
            g(x) &= 2\textcolor{red}{f(3x)} + 8 \\
             &= 2\textcolor{red}{ \frac{12x + 3}{3x-2}} + 8 \tag{substitute $f(3x)$} \\
             &=  \frac{24x + 6}{3x - 2} + 8 \tag{multiply fractions, noting $2 = \frac{2}{1}$}
        \end{align*}

        Then we can substitute $3$ into the above equation:

        \begin{align*}
            g(3) &= \frac{24 \times 3 + 6}{3 \times 3 - 2} + 8 \\
             &= \frac{78}{7} + 8 \\
             &= \frac{78 + 56}{7} \\
             &= \frac{134}{7}
        \end{align*}

        Therefore the point $P$ is mapped to $(3, g(3))$ which is the point $(3,\frac{134}{7})$.

        \item Show $ff^{-1}(x) = x$.
        This is an exercise in function composition, and we will follow the same style of approach as in e.g. Example \ref{ex:simple-composition-different-functions}.

        First write 
        \[ ff^{-1}(x) =  f(\textcolor{red}{f^{-1}(x)})\]
        Remember that 
        \[ f(x) = \frac{4x + 3}{x-2} \]
        And that we found 
        \[ f^{-1}(x) = \frac{2x+3}{x-4}\]
        Therefore 
        \begin{align*}
            f(\textcolor{red}{f^{-1}(x)}) &= \frac{4\textcolor{red}{f^{-1}(x)} + 3 }{\textcolor{red}{f^{-1}(x)} - 2} \\
                &= \frac{4\textcolor{red}{\frac{2x+3}{x-4}} + 3 }{\textcolor{red}{\frac{2x+3}{x-4}} - 2} \\ 
                &= \frac{x-4}{x-4} \times \frac{4\times \frac{2x+3}{x-4} + 3 }{\frac{2x+3}{x-4} - 2} \tag{multiply by $1$ to remove nested fractions} \\
                &= \frac{(x-4)\left(4\times \frac{2x+3}{x-4} + 3 \right)}{(x-4 )\left(\frac{2x+3}{x-4} - 2\right)} \tag{multiply fractions} \\
                &= \frac{4(3 + 2x) + 3(x-4)}{3+2x - 2(x-4)} \tag{expand} \\
                &= \frac{12 + 8x + 3x - 12}{3 + 2x - 2x + 8} \tag{expand more} \\
                &= \frac{11x}{11} \tag{collect terms}\\
                &= x
        \end{align*}
        As required, therefore $ff^{-1}(x) = x$.
    \end{enumerate}
    \qed
\end{solution}



\subsection{Domain and Range}

My school teacher said that he remembered which is range and what is domain by thinking of a lion's domain, and somehow he linked this to the inputs for a function. 
I never understood quite what he meant, but it stuck with me.

The domain of a function is the set of all input values for a function, and the range is the set of all possible outputs of that function, over the domain.

For example, consider the following function.
\[ f(x) = \frac{1}{x^2}\]

This function defined everywhere except for $x=0$, where we would be dividing by $0$. Therefore its (maximal) domain is every number except for $0$. 
On this domain, we can see that it is always positive since the denominator is always positive, and it gets closer to $0$ as $x$ gets larger (or more negative). 
We can also observe that as $x$ gets close to $0$, the function outputs larger and larger numbers. Because of this, you might realise\footnote{You can prove this by setting $f(x) = a$ where $a>0$, and showing you can solve for $x$.} 
that it can reach every number more than $0$. Therefore, we have found the domain and range of $f(x)$.

The domain is $x<0$ or $x>0$, and the range is $0 < f(x)$.

\begin{example}
    Consider $f(x) = 2x$, for $-2\leq x < 1$. Find the range for this function, making sure to only allow the $x$ range in the question.
\end{example}
\begin{solution}
    We know that this function is a straight line that is increasing, so therefore the range has to between whatever the leftmost point and the rightmost point is. The place to be careful is that 
    we are told that $-2\leq x < 1$, therefore it does not quite reach the rightmost point.
     
    We can find the range therefore by plugging in:

    \begin{align*}
        f(-2) &= 2\times(-2) = -4\\
        f(1) &= 2\times(1) = 2
    \end{align*}

    Therefore, when $-2\leq x < 1$, we have that $-4 \leq f(x) < 2$ is the range. This completes the question.
    \qed
\end{solution}

\subsubsection{The Modulus Function}
The modulus function is defined in two cases:

\begin{enumerate}
    \item If $x\geq0$, then $| x |= x$
    \item If $x<0$, then $| x |= -x$
\end{enumerate}

This has the effect of ``removing-the-minus-sign'' from the number.

\begin{example}
    \begin{align*}
        |3| &= 3\\
        |-3| &= 3\\
        |-20| &= 20\\
        |0| &= 0\\
        |5| &= 5
    \end{align*}
\end{example}

We are interested in how this can interact with a function $f(x)$, seen in Figure \ref{fig:function-for-modulus}.
\begin{figure}[H]
    \centering
    \incfig[0.75]{function for modulus}
    \caption{An example function}
    \label{fig:function-for-modulus}
\end{figure}

What would $|f(x)|$ look like? Since we know that the modulus function makes negative numbers positive, we care about when $f(x)$ is negative or positive.

Since $f(x)$ is on the $y$ axis, we know whether it is negative or positive based on whether it is below or above the $x$ axis. To plot the modulus of the function, 
observe that this is the same as reflecting all parts below the $x$ axis to be above it, and leaving the parts above the $x$ axis the same. 

\begin{figure}[H]
    \centering
    \incfig[0.75]{mod function for modulus}
    \caption{An example function, with the modulus of the function in red.}
    \label{fig:mod-function-for-modulus}
\end{figure}

Now what might $f(|x|)$ look like instead? Since we know that the modulus function makes negative numbers positive, we care about when $x$ is negative or positive.

Since $x$ is on the $x$ axis (of course it is!), we know whether it is negative or positive based on whether it is to the left or the right of the $y$ axis. 
When it is on the right of the $y$ axis, $x$ is positive and therefore $f(|x|)$ is unchanged. When it is on the left of the $x$ axis, the $x$ point that we plug into $f(x)$ is shifted to be 
the positive $x$ point. This means that the $y$ position of $f(|x|)$ is given by whatever the $y$ position is by going an equal amount to the right of the $y$ axis than you are on the left of it. 
This is the same as simply reflecting the right side of the $y$ axis to the left side. This is shown in Figure \ref{fig:modx-function-for-modulus}.

\begin{figure}[H]
    \centering
    \incfig[0.75]{modx function for modulus}
    \caption{An example function, with the function using the modulus of the input in blue.}
    \label{fig:modx-function-for-modulus}
\end{figure}

Now what about a combination of the above? What would $-|f(|x|)|$ look like? We can split this up into steps, as in when inputting an $x$, what steps happen to $x$? 
This works inside-out in order:

\begin{align*}
    x &\mapsto |x| \tag{taking the modulus}\\
    &\mapsto f(|x|) \tag{applying $f$}\\
    &\mapsto |f(|x|)| \tag{taking the modulus}\\
    &\mapsto -|f(|x|)| \tag{taking the minus sign}
\end{align*}

Now we can sketch each one in order, since each step individually we have done above. For example, we can 
start with $f(|x|)$, which we sketched in Figure \ref{fig:modx-function-for-modulus}. From here, the next step is to take the modulus.
From our discussion leading up to Figure \ref{fig:mod-function-for-modulus}, we know that this amounts to reflecting everything below the $y$ axis to above the $y$ axis.
We will do the same thing, but to $f(|x|)$. This is sketched in Figure \ref{fig:modmodx-function-for-modulus}.

\begin{figure}[H]
    \centering
    \incfig[0.75]{modmodx function for modulus}
    \caption{An example function, with the function using the modulus of the input, with the modulus applied in red}
    \label{fig:modmodx-function-for-modulus}
\end{figure}

And for the last step, we can take the minus sign of the red line sketched for  $|f(|x|)|$. This reflect every number about the $x$ axis. This gives us Figure \ref{fig:minusmodmodx-function-for-modulus}.
\begin{figure}[H]
    \centering
    \incfig[0.75]{minusmodmodx function for modulus}
    \caption{An example function, with the function using the modulus of the input, with the modulus applied, with a negative sign in red}
    \label{fig:minusmodmodx-function-for-modulus}
\end{figure}

The final result would have been hard to spot directly from Figure \ref{fig:function-for-modulus} without breaking it down step-by-step.

We are often interested in how the range of a function is affected by the modulus.

\begin{example}
    Given that $f(x)$ has range $-16 < f(x) \leq 5$ and reaches every point inbetween, what range does $|f(x)|$ have?
\end{example}
\begin{solution}
    Since the range means that the $y$ values are between $-16$ and $5$ (but not equal to $-16$ in this case), and we know the 
    modulus turns negative values into positive values, we therefore know that $-16$ will transform into $16$. Since this is more than $5$, this is the new top-end of the range. 
    For the lower end, we are told that $f$ reaches every point inbetween. Therefore the lowest point the modulus can reach is $0$.

    This gives us the range $0 \leq |f(x)| < 16$, since the original function did not quite reach $-16$, the modulus can't reach $16$ either. This completes the question.
    \qed
\end{solution}
The key point here is that the modulus does not ``remove'' values; it makes them positive. One way to think about the previous question is that 
$-16$ is ``further negative'' than $5$ is positive, the new upper end of the range must go to the one that is further out, in this case $|-16|$.


\begin{example}
    Let $f(x) = 2x - 1$, for $-2 \leq x \leq 1$.

    \begin{enumerate}
        \item What is the range of $f(x)$?
        \item What is the range of $|f(x)|$?
        \item What is the range of $f(|x|)$ if we are given that $f$ is not defined for inputs not between $-2$ and $-1$? What is the range if we allowed $-2 \leq x \leq 1$ for $f(|x|) = 2|x| - 1$?\footnote{The difference here is whether we impose that $f$ itself is only defined on inputs between $-2$ and $1$, or if we consider $f(|x|)$ on all inputs between $-2$ and $1$, which actually implies that we are allowed to use $f$ between the points $0$ and $2$, since we are taking the modulus first before plugging into $f$. This is an important clarification since $2$ is outside our original allowed range of $x$ values.}
    \end{enumerate}
\end{example}

\begin{solution}

    \begin{enumerate}
        \item 
    We know that this function is a straight line that is increasing, so therefore the range has to between whatever the leftmost point and the rightmost point is. The place to be careful is that 
    we are told that $-2\leq x \leq 1$.
     
    We can find the range therefore by plugging in:

    \begin{align*}
        f(-2) &= 2\times(-2) - 1 = -4 -1 = -5\\
        f(1) &= 2\times(1) -1= 2 -1 = 1
    \end{align*}

    Therefore, when $-2\leq x < 1$, we have that $-5 \leq f(x) \leq 1$ is the range. This completes the first part of the question.
        \item
    For this part, it is useful to sketch the graph, Figure \ref{fig:line-for-modulus}

    \begin{figure}[H]
    \centering
    \incfig[0.75]{line for modulus}
    \caption{(Sketch) The function $2x - 1$, with the modulus of the function in red, for $-2 \leq x \leq 1$.}
    \label{fig:line-for-modulus}
    \end{figure}

    From here we can see that the new minimum of the function is $0$, and the new maximum is the old minimum, which in this case is $|-5|= 5$.

    This gives us the range $0 \leq |f(x)| \leq 5$, which completes this part of the question.

    \item 
    For this part, it is again useful to sketch the graph, Figure \ref{fig:modxline-for-modulus}

    \begin{figure}[H]
    \centering
    \incfig[0.75]{modxline for modulus}
    \caption{(Sketch) The function $2x - 1$, with the function using the modulus of the input in red, for $-2 \leq x \leq 1$.}
    \label{fig:modxline-for-modulus}
    \end{figure}
    
    Notice that this the function cuts off on the left side. This is because we are only allowed to use $f(x)$ when $x$ is between $-2$ and $1$ from the question. But if $x$ is less than $-1$, then $|x|$ is more than $1$, which is not allowed to be input into the function.
    
    From here we can see that the new minimum of the function is where the old function intersected the $y$ axis, at $x=0$, and the new maximum is the maximum of the right side of $f(x)$, which in this case is at $x=1$ or $x=-1$.

    At $x=0$, 
    \[f(|0|) = 2\times 0 - 1 = -1\] 
    This gives the minimum as $-1$. at $x=1$, 
    \[ f(|1|) = 2 \times 1 - 1 = 1\]

    This gives us the range $-1 \leq |f(x)| \leq 1$, which completes this part of the question.
    
    If we allow $-2 \leq x \leq 1$ as in the second half of this question, the red line would extend backwards to $x=-2$, where $f(|x|) = 2 \times |-2| - 1 = 3$, and therefore the 
    range in this case is $-1 \leq |f(x)| \leq 3$.
    
    This completes the full question.
\end{enumerate}
    \qed
\end{solution}

